% paper_21_geometric_vacuum_synthesis.tex — GeoVac Paper 21
% The Geometric Vacuum: Quantum Structure from Compact Manifold Topology
% Status: Draft, April 7, 2026
% Category: Synthesis (papers/synthesis/)
\documentclass[aps,pra,twocolumn,superscriptaddress,longbibliography]{revtex4-2}

\usepackage{amsmath,amssymb,amsthm,graphicx,xcolor,bm,braket,dcolumn,booktabs}
\usepackage{hyperref}

\newtheorem{theorem}{Theorem}
\newtheorem{observation}{Observation}
\newtheorem{conjecture}{Conjecture}

\begin{document}

\title{The Geometric Vacuum:\\
Quantum Structure from Compact Manifold Topology}
\thanks{Paper 21 in the Geometric Vacuum series}

\author{J.~Loutey}
\affiliation{Independent Researcher, Kent, Washington}

\date{April 7, 2026}

%% ========================================================================
%% ABSTRACT
%% ========================================================================

\begin{abstract}
The non-relativistic quantum mechanics of Coulomb systems can be
reformulated as spectral geometry on compact manifolds.  The angular
momentum quantum numbers $(l, m)$ and their degeneracies $2l+1$ emerge
from a geometric packing construction on isotropic shells, independent
of the Schr\"odinger equation.  The hydrogen energy spectrum, the
$1/r$ Coulomb potential, and the $n^2$ degeneracy all follow from the
Laplace--Beltrami operator on the unit three-sphere $S^3$, with physical
scales entering exclusively through stereographic projection---as
demonstrated by Fock in 1935 and verified here through 18 independent
symbolic proofs.  For $N$-electron atoms, the framework extends to
$S^{3N-1}$, where the permutation group $S_N$ acting on the
hyperspherical angular basis produces chemical periodicity as a
representation-theoretic consequence.  A new sparsity mechanism from
angular momentum recoupling---6$j$ coefficients producing additional
zeros beyond Gaunt selection rules---reduces electron repulsion integral
density by 18--26\%.  The exchange constant taxonomy classifies
transcendental content (where $\pi$ and other transcendentals enter the
otherwise rational graph spectrum) as instances of the Weyl--Selberg
spectral-geometric conversion.  The Hopf fibration
$S^1 \to S^3 \to S^2$ of the electron's proven state manifold yields
the fine structure constant $\alpha^{-1} = 137.035987$ from spectral
invariants of the bundle components, with zero free parameters and
$p$-value $5.2 \times 10^{-9}$; the combination rule remains conjectural.
These results are mathematically exact for single-center systems and
suggest a concrete research program: high-precision atomic spectra on
$S^{3N-1}$, compared with experimental fine structure data, can test the
$\alpha$ conjecture computationally.
\end{abstract}

\maketitle

%% ========================================================================
%% I. INTRODUCTION
%% ========================================================================

\section{Introduction}
\label{sec:introduction}

The Schr\"odinger equation for hydrogen presupposes continuous flat
space, a radial Coulomb potential $V(r) = -Z/r$, and quantized energy
levels $E_n = -Z^2/(2n^2)$ as boundary-condition eigenvalues.  Fock
showed in 1935~\cite{Fock1935} that these are not independent
assumptions: the hydrogen problem on the unit three-sphere $S^3$ is a
free particle, with the $1/r$ potential emerging from stereographic
projection and the Rydberg spectrum from the energy-shell constraint
$p_0^2 = -2E$.  Bargmann~\cite{Bargmann1936} identified the underlying
SO(4) symmetry, and Barut and Kleinert~\cite{BarutKleinert1967}
extended the dynamical algebra to SO(4,2).

These foundational results are well known.  What is new in the present
work is their extension into a complete computational framework:
\begin{enumerate}
\item A geometric packing construction derives the angular quantum
  numbers $(l, m)$ and their degeneracies $2l+1$ from two axioms
  (binary distinguishability and geometric indifference), independent
  of operators or dynamics (Paper~0~\cite{Paper0}).
\item The discrete graph Laplacian on the resulting lattice converges
  to the Laplace--Beltrami operator on the unit $S^3$, verified
  through 18 independent symbolic proofs (Paper~7~\cite{Paper7}).
\item The framework extends to $N$-electron atoms on $S^{3N-1}$, where
  the permutation group $S_N$ produces chemical periodicity as
  representation theory (Papers~13, 16~\cite{Paper13,Paper16}).
\item A new sparsity mechanism from angular momentum
  recoupling---6$j$ coefficients producing additional zeros beyond
  Gaunt selection rules---yields 18--26\% lower electron repulsion
  integral (ERI) density in coupled bases.
\item The composed molecular architecture assigns each electron group
  its own natural coordinate system, producing structurally sparse qubit
  Hamiltonians: $N_{\text{Pauli}} \approx 11.11 \times Q$ across a
  library of 30 molecules spanning three rows of the periodic table
  ($Z = 1$--36), with open-source ecosystem export to OpenFermion,
  Qiskit, and PennyLane
  (Papers~14, 17, 19, 20~\cite{Paper14,Paper17,Paper19,Paper20}).
\item The Hopf fibration $S^1 \to S^3 \to S^2$ suggests a topological
  origin for the electromagnetic coupling constant $\alpha$
  (Paper~2~\cite{Paper2}, conjectural).
\end{enumerate}

\noindent\textit{Scope.}  The reformulation is mathematically exact for
single-center Coulomb systems (atoms), where the $S^{3N-1}$ structure
provides 18 symbolic proofs, 0.004\% He ground-state accuracy, and a
new 6$j$ recoupling sparsity mechanism.  For molecules, the composed
architecture factorizes the $N$-body problem into coupled natural
geometries, producing practical qubit Hamiltonians with $O(Q^{2.5})$
Pauli scaling and characterized FCI accuracy (0.20\% energy error for
LiH at fixed geometry via the balanced coupled Hamiltonian).  This paper
covers both: the single-center mathematical foundations
(Secs.~\ref{sec:kinematics}--\ref{sec:exchange}) and the molecular
computational pipeline that builds on them (Sec.~\ref{sec:pipeline}).

%% ========================================================================
%% II. KINEMATICS: QUANTUM NUMBERS FROM GEOMETRY
%% ========================================================================

\section{Kinematics: Quantum Numbers from Geometry}
\label{sec:kinematics}

The angular structure of quantum mechanics---the quantum numbers
$(l, m)$, the degeneracy $2l + 1$ per shell, and the cumulative count
$n^2$---can be derived from two axioms without reference to the
Schr\"odinger equation, angular momentum operators, or any specific
physical system.

\subsection{Two axioms}

\begin{enumerate}
\item \textbf{Binary distinguishability.}  A minimum of two
  distinguishable states is required to establish a nonzero fundamental
  scale $d_0$.  A single point cannot define a distance; two points
  create the first bit of spatial information.
\item \textbf{Geometric indifference.}  In the absence of external
  constraints or privileged directions, the distribution of states
  must maximize entropy.  By the principle of indifference, this
  necessitates isotropy: states populate isotropic shells (circles)
  rather than anisotropic tilings (squares, hexagons), which would
  imply privileged directions not present in the axioms.
\end{enumerate}

\subsection{Packing construction}

The two-point initialization fixes the fundamental area per state
$\sigma_0 = \pi d_0^2 / 2$.  States are distributed on linearly
spaced shells with radius $r_k = k \cdot d_0$, $k = 1, 2, 3, \ldots$
The annular area of shell $k$ is
\begin{equation}
A_k = \pi d_0^2 \bigl[ k^2 - (k-1)^2 \bigr] = \pi d_0^2 (2k - 1),
\label{eq:annular_area}
\end{equation}
and the angular capacity per shell is
\begin{equation}
N_k^{\text{angular}} = \frac{A_k}{\sigma_0} = 2(2k - 1).
\label{eq:shell_capacity}
\end{equation}
The factor $(2k-1)$ is the dimension of the spin-$(k-1)$ irreducible
representation of SO(3).  With the identification $l = k - 1$, each
shell carries exactly $2l + 1$ angular states---the spherical harmonic
basis $Y_{lm}$ with $m = -l, \ldots, +l$.  The global factor of 2
arises from orientability of the compactified packing manifold on $S^2$.

The cumulative state count through shell $n$ is
\begin{equation}
\sum_{k=1}^{n} (2k - 1) = n^2,
\label{eq:cumulative}
\end{equation}
reproducing the $n^2$ orbital count (or $2n^2$ including the spin
factor).

\subsection{What the construction establishes}

The packing construction derives the \textit{kinematic} structure of
quantum mechanics---the state-counting aspect of angular momentum---from
pure geometry.  It produces the quantum number labeling $(l, m)$, the
per-shell degeneracy $2l+1$, and the cumulative count $n^2$.  It does
not produce energy eigenvalues, wavefunctions, or dynamics; those
require the \textit{dynamic} structure provided by $S^3$ (Sec.~\ref{sec:dynamics}).

%% ========================================================================
%% III. DYNAMICS: THE SCHRÖDINGER EQUATION FROM S³
%% ========================================================================

\section{Dynamics: The Schr\"odinger Equation from $S^3$}
\label{sec:dynamics}

\subsection{The paraboloid lattice}

The quantum numbers $(n, l, m)$ from the packing construction label
nodes of a discrete graph---the paraboloid lattice.  Edges are generated
by the ladder operators $L_\pm$ (angular momentum raising/lowering) and
$T_\pm$ (radial transitions), producing a sparse adjacency structure.
The graph Hamiltonian is
\begin{equation}
H = \kappa \cdot \mathcal{L},
\label{eq:graph_hamiltonian}
\end{equation}
where $\mathcal{L} = D - A$ is the graph Laplacian ($D$ = degree
matrix, $A$ = adjacency matrix) and $\kappa = -1/16$ is the universal
kinetic scale factor.

\subsection{Continuum limit: the unit $S^3$}

In the continuum limit ($n_{\max} \to \infty$), the graph Laplacian
converges to the Laplace--Beltrami operator $\Delta_{S^3}$ on the unit
three-sphere.  The eigenfunctions are Gegenbauer polynomials
$C_{n-1}^1(\cos\chi)$ in the $l = 0$ sector, where $\chi$ is the
hyperspherical angle, and the eigenvalues are pure dimensionless
integers:
\begin{equation}
\Delta_{S^3}\, C_{n-1}^1(\cos\chi) = -(n^2 - 1)\, C_{n-1}^1(\cos\chi).
\label{eq:s3_eigenvalues}
\end{equation}
These eigenvalues $\lambda_n = -(n^2 - 1)$ are topological invariants
of the unit sphere, carrying no physical dimensions.

\subsection{Stereographic projection: Fock's map}

Fock's stereographic projection~\cite{Fock1935} maps the unit $S^3$ to
flat momentum space $\mathbb{R}^3$.  A point
$\mathbf{n} = (n_1, n_2, n_3, n_4) \in S^3$ maps to momentum
$\mathbf{p} = (p_1, p_2, p_3) \in \mathbb{R}^3$ via
\begin{equation}
n_i = \frac{2p_0\, p_i}{p^2 + p_0^2}, \quad
n_4 = \frac{p^2 - p_0^2}{p^2 + p_0^2},
\label{eq:stereographic}
\end{equation}
with conformal factor
\begin{equation}
\Omega(\mathbf{p}) = \frac{2p_0}{p^2 + p_0^2}.
\label{eq:conformal_factor}
\end{equation}
The unit sphere constraint $\sum_{i=1}^{4} n_i^2 = 1$ holds for
\textit{all} $p_0 > 0$, confirming scale invariance: the projection
always produces the unit $S^3$ regardless of the choice of $p_0$.

\subsection{Emergence of the Coulomb potential}

The chordal distance identity
\begin{equation}
|\mathbf{n}(\mathbf{p}) - \mathbf{n}(\mathbf{q})|^2
= \Omega(\mathbf{p})\, \Omega(\mathbf{q})\, |\mathbf{p} - \mathbf{q}|^2
\label{eq:chordal}
\end{equation}
relates the Euclidean distance on $S^3$ (left side) to the flat-space
distance (right side), weighted by conformal factors.  The free-particle
Green's function on $S^3$, which depends on the chordal distance
$|\mathbf{n} - \mathbf{n}'|^{-2}$, transforms under stereographic
projection into a kernel proportional to
$1/|\mathbf{p} - \mathbf{p}'|^2$.  Upon Fourier transformation to
position space, this becomes the $1/r$ Coulomb potential.

The Coulomb potential is therefore not an input force law---it is the
coordinate distortion introduced by stereographic projection.

\subsection{Energy-shell constraint and the Rydberg spectrum}

The parameter $p_0$ is fixed by the energy-shell constraint
\begin{equation}
p_0^2 = -2E,
\label{eq:energy_shell}
\end{equation}
which converts the dimensionless integer eigenvalues into physical
energies.  For the $n$-th level, $p_0 = 1/n$ (in atomic units), and
\begin{equation}
E_n = -\frac{p_0^2}{2} = -\frac{1}{2n^2}.
\label{eq:rydberg}
\end{equation}
The complete chain is:
\begin{equation}
\text{Graph} \;\xrightarrow{\;\text{continuum}\;}\; S^3
\;\xrightarrow{\;\text{stereo.}\;}\; \mathbb{R}^3
\;\xrightarrow{\;p_0^2 = -2E\;}\; E_n = -\frac{1}{2n^2}.
\label{eq:chain}
\end{equation}

\subsection{Symbolic verification}

Every algebraic step in this chain has been verified through 18
independent symbolic proofs~\cite{Paper7}, implemented in the
\texttt{sympy} computer algebra system and collected in the test suite
\texttt{test\_fock\_projection.py} and \texttt{test\_fock\_laplacian.py}.
These verify:
\begin{itemize}
\item Module I (Tests 1--10): Stereographic projection geometry---unit
  sphere constraint, conformal factor identities, chordal distance
  relation, volume Jacobian, inverse projection round-trip.
\item Module II (Tests 11--18): Conformal Laplacian and
  eigenvalues---connection term, zero-mode annihilation
  ($\Delta_{S^3}(1) = 0$), eigenvalue verification at $n = 2$ and
  $n = 3$, rational decomposition of conformally weighted Laplacian,
  eigenvalue-to-energy mapping.
\end{itemize}
All 18 tests pass with symbolic (exact) agreement, constituting a
machine-verified proof of the graph-to-Schr\"odinger equivalence.

%% ========================================================================
%% IV. THE MULTI-ELECTRON GENERALIZATION
%% ========================================================================

\section{The Multi-Electron Generalization}
\label{sec:multielectron}

\subsection{$N$ electrons on $S^{3N-1}$}

For an $N$-electron atom, the $3N$ spatial coordinates are collected
into a hyperradius $R = \sqrt{\sum_i r_i^2}$ and $3N-1$ angular
coordinates on $S^{3N-1}$.  The kinetic energy operator separates as
\begin{equation}
T = -\frac{1}{2}\left(\frac{\partial^2}{\partial R^2}
+ \frac{3N-1}{R}\frac{\partial}{\partial R}
- \frac{\Lambda^2}{R^2}\right),
\label{eq:kinetic_N}
\end{equation}
where $\Lambda^2$ is the grand angular momentum operator (the
Laplace--Beltrami operator on $S^{3N-1}$).  Its eigenvalues are the
SO($3N$) Casimir invariants:
\begin{equation}
\Lambda^2 Y_\nu = \nu(\nu + 3N - 2)\, Y_\nu,
\label{eq:casimir_3N}
\end{equation}
with $\nu = 0, 1, 2, \ldots$ and $Y_\nu$ the hyperspherical harmonics.

At hyperradius $R = 0$, the system is a free particle on the unit
$S^{3N-1}$.  The total Coulomb interaction enters as a geometric
perturbation through the charge function $C(\hat{\Omega})/R$, where
$C$ encodes the nuclear attraction and electron-electron repulsion in
angular coordinates.  In the adiabatic hyperspherical
method~\cite{Macek1968,Lin1995}, the angular eigenvalue problem is
solved parametrically at each $R$, yielding adiabatic eigenvalues
$\mu_\nu(R)$ that define effective radial potentials
$V_{\text{eff}}(R) = \mu_\nu(R)/R^2 + (3N-1)(3N-3)/(8R^2)$.

For helium ($N = 2$, $S^5$), the single-channel adiabatic approximation
yields a ground-state energy of $-2.9052$~Ha, a 0.05\% error relative
to the exact nonrelativistic value $-2.903724$~Ha~\cite{Paper13}.

\subsection{Chemical periodicity from $S_N$ representation theory}

The total wavefunction must be antisymmetric under electron permutation.
Writing $|\Psi\rangle = |\text{spatial}\rangle \otimes |\text{spin}\rangle$
and applying Schur--Weyl duality, the spatial symmetry is determined by
the conjugate of the spin Young diagram.

\begin{theorem}[Universal $\lambda_1$~\cite{Paper16}]
For any $N$-electron atom with ground-state spin $S < N/2$, the spatial
Young diagram has first-column length $\lambda_1 = 2$, independent of
$Z$, $S$, or electron configuration.
\end{theorem}

\noindent\textit{Proof.}  The spin diagram has two rows
$[n_\uparrow, n_\downarrow]$ with $n_\downarrow \geq 1$ (since
$S < N/2$).  Its conjugate has $n_\downarrow$ columns of height~2 and
$2S$ columns of height~1.  The first-column length is~2.  \qed

This fixes the minimum grand angular momentum quantum number at
$\nu = N - \lambda_1 = N - 2$, giving a universal free eigenvalue:
\begin{equation}
\boxed{\mu_{\text{free}} = \frac{(N-2)(4N-4)}{2} = 2(N-2)^2.}
\label{eq:mu_free}
\end{equation}
This ``Pauli centrifugal cost'' depends only on the electron count $N$,
not on the nuclear charge $Z$, spin state $S$, or orbital
configuration---it is purely topological.

The periodic table structure follows from the sequence of $S_N$ spatial
irreducible representations as shells fill.  Five structure types emerge:
\begin{description}
\item[Type A] ($N = 1$): Hydrogen.  $\nu = 0$, no antisymmetry
  constraint.
\item[Type B] (Noble gases): Closed-shell, democratic Young diagram
  $[2^{N/2}]$.  Maximal Pauli centrifugal barrier.
\item[Type C] (Alkali metals): One electron in new shell above closed
  core.  Hierarchical diagram $[2^{(N-1)/2}, 1]$.
\item[Type D] (Alkaline earths): Two electrons above closed core.
  Same diagram shape as Type~B but different occupancy pattern.
\item[Type E] (Open-shell): Partially filled $p$- and $d$-subshells.
  Varying spin and diagram shapes.
\end{description}

The periodic law is the group-theoretic statement: \textit{chemical
properties of elements are periodic functions of atomic number because
the sequence of $S_N$ spatial irreps repeats whenever a new principal
shell begins filling.}

\subsection{Angular momentum recoupling sparsity}

For $N$-electron systems, organizing the $S^{3N-1}$ angular basis
hierarchically via pairwise (H-set) coupling---where electrons are
grouped into pairs with definite intermediate angular momenta
$l_{12}, l_{34}, \ldots$---produces a new sparsity mechanism.

When a two-body operator (e.g., the electron-electron repulsion
$V_{ee}(1,2)$) acts on a pair, it preserves the angular momenta of
spectator pairs.  The matrix element of $V_{ee}$ between coupled states
involves Wigner 6$j$ symbols from recoupling:
\begin{equation}
\begin{Bmatrix} l_1 & l_2 & l_{12} \\ l_{12}' & l_1' & k \end{Bmatrix},
\label{eq:6j}
\end{equation}
where $k$ is the multipole order from the Gaunt integral.  These 6$j$
coefficients vanish for certain combinations of angular momenta,
producing \textit{additional zeros} beyond the bare Gaunt selection
rules ($|l - l'| \leq k \leq l + l'$, $l + k + l' = \text{even}$).

The resulting ERI density reduction is quantified in
Table~\ref{tab:eri_density}.

\begin{table}[t]
\caption{ERI density (fraction of nonzero electron repulsion integrals)
in uncoupled vs.\ H-set coupled bases for four-electron systems.}
\label{tab:eri_density}
\begin{ruledtabular}
\begin{tabular}{ccccc}
$l_{\max}$ & Uncoupled & H-set & Reduction \\
\hline
1 & 12.5\% & 9.2\% & 26\% \\
2 & 5.1\% & 4.2\% & 18\% \\
\end{tabular}
\end{ruledtabular}
\end{table}

This sparsity mechanism is basis-agnostic: it applies to any encoding
with angular momentum pair coupling and follows from the representation
theory of SO(3), not from any particular choice of radial basis or
physical system.

\paragraph{Three-layer structure of the He ground state (v2.6.0).}
The graph-native FCI (Track~DI Sprints~2--3D) makes the exchange constant
structure fully explicit.  Using the graph Laplacian $\kappa\mathcal{L}$
as the one-body operator and exact rational Slater integrals as the
two-body operator, the He ground state is computed with zero free
parameters.  At $n_{\max} = 7$ (1,218~configurations): 0.19\% error.
Replacing $\kappa\mathcal{L}$ with the diagonal hydrogenic operator
degrades to 1.41\% (variational~$k$) or 1.85\% (fixed $k = Z$),
demonstrating that the graph topology (off-diagonal $T_\pm$ couplings)
accounts for $\sim$86\% of the CI correlation energy.

FCI basis invariance was verified (Sprint~3D): transforming $V_{ee}$
to the graph eigenbasis gives identical energies ($< 3 \times 10^{-15}$~Ha),
confirming the $\sim$0.14\% convergence floor is irreducible
transcendental content from the electron-electron cusp at
$\alpha = \pi/4$ on~$S^5$ (embedding exchange constant, Paper~18),
not basis mismatch between the graph and hydrogenic operators.

%% ========================================================================
%% V. THE EXCHANGE CONSTANT TAXONOMY
%% ========================================================================

\section{Where Continuous Physics Enters}
\label{sec:exchange}

\subsection{The $\pi$-free graph}

The graph Laplacian on the paraboloid lattice is $\pi$-free: all
eigenvalues, degeneracies, and coupling coefficients are integers or
rationals.  Transcendental numbers---$\pi$, exponential integrals,
spectral zeta values---enter exclusively when projecting the discrete
graph onto continuous manifolds for comparison with experiment or
computational convenience.

This observation motivates a classification of exactly
\textit{where} continuous physics enters the otherwise discrete
framework.

\subsection{Four types of exchange constant}

Each projection from the discrete algebraic structure to a continuous
coordinate system introduces a characteristic transcendental constant,
classified by what determines it~\cite{Paper18}:

\begin{description}
\item[Intrinsic:] Determined by dimension alone.  Example: $\pi$ for
  $S^1 \to \mathbb{R}$.

\item[Calibration:] Determined by graph connectivity.  Example: the
  kinetic scale factor $\kappa = -1/16$, which maps the graph
  eigenvalue $\lambda_{\max} = 8$ to the hydrogen ground state
  $E_H = -1/2$~Ha.

\item[Embedding:] Seeded by coordinate singularities.  Example: the
  Stieltjes seed $e^a E_1(a)$ that enters when embedding the Laguerre
  spectral basis into prolate spheroidal coordinates (Paper~11), and
  the $1/r_{12}$ cusp---an embedding exchange constant that cannot be
  absorbed into the $S^5$ graph
  (Paper~15~\cite{Paper15}).

\item[Flow:] Parameter-dependent algebraic functions with
  transcendental coefficients.  Example: the adiabatic eigenvalue
  $\mu(R)$, which satisfies a polynomial equation
  $P(R, \mu) = 0$ at each hyperradius but whose $R$-dependence is
  transcendental (Paper~13).
\end{description}

\noindent These form a dependency hierarchy: intrinsic $\to$ calibration
$\to$ embedding $\to$ flow, with each level incorporating the
transcendental content of the previous levels.

\subsection{Weyl--Selberg identification}

The exchange constants are instances of a known mathematical
structure.  Weyl's law relates the eigenvalue counting function
$N(\Lambda)$ of a compact Riemannian manifold $\mathcal{M}$ to its
volume:
\begin{equation}
N(\Lambda) \sim \frac{\mathrm{Vol}(\mathcal{M})}{(2\pi)^d}\,
\Lambda^{d/2},
\label{eq:weyl}
\end{equation}
providing the precise conversion between discrete eigenvalue sums and
continuous manifold volumes.  The Selberg trace formula generalizes
this to an exact identity between spectral data and geometric data
(lengths of closed geodesics) on hyperbolic manifolds.

The GeoVac exchange constants---$\kappa$, $e^a E_1(a)$,
$\mu(R)$---are specific instances of this spectral-geometric
conversion for the conformal and hyperspherical hierarchies.  The
exchange constant taxonomy tells us exactly where continuous physics
enters the discrete framework, and why it takes the specific
transcendental form it does: the graph is $\pi$-free; $\pi$ enters
when you project it.

\subsection{Observable classification}

The exchange constant hierarchy induces a classification of physical
observables by their transcendental content~\cite{Paper18}:

\begin{description}
\item[Class S (Spectral):] Energy eigenvalues, degeneracies, selection
  rules.  Require graph eigenvalues and $\kappa$ only.  \textit{No
  transcendental content} beyond the calibration constant.

\item[Class P (Single-particle spatial):] Electron density $\rho(r)$,
  radial moments $\langle r^k \rangle$, momentum distributions.
  Require the $S^3 \to \mathbb{R}^3$ conformal projection, introducing
  $\pi$ through $\mathrm{Vol}(S^3) = 2\pi^2$.

\item[Class C (Multi-particle correlated):] Correlation energies,
  adiabatic potential energy surfaces, multi-electron expectation
  values.  Require embedding and flow constants ($e^a E_1(a)$,
  $\mu(R)$, PK pseudopotential).
\end{description}

\noindent
The He ground state provides a concrete census of all four exchange
constant types (Paper~18, Note added v2.6.0): intrinsic ($\pi$ from
$\mathrm{Vol}(S^5)$), calibration ($\kappa = -1/16$), embedding
(e-e cusp at $\alpha = \pi/4$), and flow ($\mu(R)$, bypassed by the
2D variational solver).  The three-layer decomposition---rational
Slater integrals, algebraic optimal exponents, transcendental
cusp---demonstrates that each layer corresponds to a distinct exchange
constant class.

%% ========================================================================
%% VI. THE HOPF FIBRATION AND ELECTROMAGNETIC COUPLING
%% ========================================================================

\section{The Hopf Fibration and Electromagnetic Coupling}
\label{sec:alpha}

The results of Secs.~\ref{sec:dynamics}--\ref{sec:exchange} establish
$S^3$ as the proven state manifold for a single bound electron.
The three-sphere possesses a unique nontrivial fiber bundle---the Hopf
fibration~\cite{Hopf1931}---which we now examine for physical content.

\textit{The results of this section are conjectural.  The numerical
agreement is statistically significant but the combination rule
(Eq.~\ref{eq:combination_rule}) has not been derived from first
principles.  We present the evidence and identify explicitly which
links in the derivation chain are established, partially established,
and open.}

\subsection{The Hopf bundle as electron-photon geometry}

The Hopf fibration
\begin{equation}
S^1 \;\longrightarrow\; S^3 \;\longrightarrow\; S^2
\label{eq:hopf}
\end{equation}
is the unique principal $U(1)$-bundle over $S^2$ with first Chern
number $c_1 = 1$.  Its three components admit a physical interpretation:
\begin{itemize}
\item $S^3$ (total space): The electron bound-state manifold, proven by
  Fock's stereographic projection (Sec.~\ref{sec:dynamics}).
\item $S^2$ (base): The photon momentum sphere, parametrizing
  propagation directions of the electromagnetic field.
\item $S^1$ (fiber): The $U(1)$ gauge phase of the electromagnetic
  potential.
\end{itemize}
Electromagnetic coupling is then the cost of projecting the electron's
$S^3$ onto the photon's $S^2$---a projection whose geometry is entirely
determined by the Hopf map.

\subsection{Spectral invariants}

Each bundle component contributes a spectral invariant.

\textit{Base $S^2$---Casimir trace.}  The SO(3) Casimir eigenvalue for
angular momentum $l$ is $l(l+1)$, with degeneracy $2l+1$.  Summing
over the hydrogen shells up to a cutoff $n_{\max}$:
\begin{equation}
B(n_{\max}) = \sum_{n=1}^{n_{\max}} \sum_{l=0}^{n-1}
(2l+1)\, l(l+1).
\label{eq:casimir_trace}
\end{equation}
At $n_{\max} = 3$: $B(3) = 0 + 0 + 6 + 0 + 6 + 30 = 42$.  The
closed-form expression is $B(m) = m(m+1)(2m+1)(m+2)(m-1)/20$.

\textit{Fiber $S^1$---spectral zeta function.}  The $S^1$ zeta function
at $s = 1$ gives
\begin{equation}
F = \zeta_R(2) = \frac{\pi^2}{6} = 1.644934\ldots
\label{eq:fiber_zeta}
\end{equation}

\textit{Total $S^3$---boundary correction.}  At cutoff $n_{\max} = 3$,
the highest eigenvalue is $|\lambda_3| = 8$ and the cumulative state
count through shell~2 is $N(2) = 5$:
\begin{equation}
\Delta = \frac{1}{|\lambda_3| \times N(2)} = \frac{1}{40} = 0.025.
\label{eq:boundary}
\end{equation}

\subsection{The selection principle}

The cutoff $n_{\max} = 3$ is not arbitrary.  The ratio $B/N$, where
$N(n_{\max}) = \sum_{n=1}^{n_{\max}} n^2$ is the total state count,
equals the dimension of $S^3$ at exactly one value:
\begin{equation}
\frac{B(n_{\max})}{N(n_{\max})} = \dim(S^3) = 3
\quad\Leftrightarrow\quad n_{\max} = 3.
\label{eq:selection}
\end{equation}
This follows from the closed-form ratio
$B(m)/N(m) = 3(m+2)(m-1)/10$; setting this equal to~3 gives
$(m+4)(m-3) = 0$, with unique positive root $m = 3$.  The proof is
algebraic and specific to $S^3$: the $SO(4) \cong (SU(2) \times
SU(2))/\mathbb{Z}_2$ product structure yields the clean
Clebsch--Gordan branching $l = 0, \ldots, n-1$ within each shell,
and $S^3 \cong SU(2)$ is the only sphere that is a Lie group.

\subsection{The combination rule and cubic equation}

The conjectural combination rule is
\begin{equation}
K = \pi\!\left(B + F - \Delta\right) = \pi\!\left(42 + \frac{\pi^2}{6}
- \frac{1}{40}\right) = 137.036064.
\label{eq:combination_rule}
\end{equation}
The self-consistency condition $\alpha^{-1} + \alpha^2 = K$ yields
the depressed cubic
\begin{equation}
\alpha^3 - K\alpha + 1 = 0,
\label{eq:cubic}
\end{equation}
whose discriminant $4K^3 - 27 > 0$ confirms three real roots.
Selecting the smallest positive root (corresponding to perturbative
QED):
\begin{equation}
\alpha^{-1} = 137.035987.
\label{eq:alpha_result}
\end{equation}
The CODATA~2018 experimental value is
$\alpha^{-1} = 137.035999\ldots$, giving a relative error of
$8.8 \times 10^{-8}$ with zero free parameters.

\subsection{Circulant structure}

The cubic~(\ref{eq:cubic}) is the characteristic polynomial of a
traceless $\mathbb{Z}_3$-symmetric circulant matrix:
\begin{equation}
M = \begin{pmatrix} 0 & b & c \\ c & 0 & b \\ b & c & 0 \end{pmatrix},
\label{eq:circulant}
\end{equation}
with $bc = K/3$ and $b^3 + c^3 = -1$.  The $U(1)$ gauge structure of
the Hopf fiber forces Hermiticity: $c = b^*$, so
$b = \sqrt{K/3}\, e^{i\theta}$ with phase determined by
$\cos(3\theta) = -1/[2(K/3)^{3/2}] \approx -0.00162$.  The physical
root corresponds to $\theta \approx 3\pi/2$ (nearly pure imaginary
coupling), encoding a democratic interaction among three bundle
components.

\subsection{Statistical significance}

To assess whether the numerical agreement is coincidental, a
combinatorial search was conducted over $1.92 \times 10^9$ formulas of
the form $\alpha^{-1} = T(A + B + C)$ (or cubic variants), where
$A, B, C$ are drawn from 743 spectral terms and $T$ from 14
transcendental prefactors~\cite{Paper2}.  Only 10 formulas achieve
precision $\leq 8.8 \times 10^{-8}$, giving
\begin{equation}
p\text{-value} = 5.2 \times 10^{-9}.
\label{eq:pvalue}
\end{equation}
This distinguishes the result from numerology but does not constitute a
derivation.

\subsection{Derivation chain}

Table~\ref{tab:chain} summarizes the status of each link.

\begin{table}[t]
\caption{Derivation chain for $\alpha$ from the Hopf bundle.}
\label{tab:chain}
\begin{ruledtabular}
\begin{tabular}{clc}
Link & Description & Status \\
\hline
1 & Hopf bundle as geometric arena & $\checkmark$ \\
2 & Spectral data $\to$ $B$, $F$, $\Delta$ & $\bullet$ \\
3 & Combination rule $K = \pi(B+F-\Delta)$ & $\times$ \\
4 & $\mathbb{Z}_3$ circulant $\to$ cubic & $\circ$ \\
5 & $\alpha$ from cubic root & $\checkmark$ \\
\end{tabular}
\end{ruledtabular}
\begin{flushleft}
$\checkmark$ = established; $\bullet$ = largely established
(selection principle $B/N = 3$ proven; gap: why these specific
spectral quantities); $\circ$ = partially established (Hermiticity
proven; gap: why $bc = K/3$); $\times$ = not established (critical
gap: no principle selects this additive combination).
\end{flushleft}
\end{table}

Link~3---the combination rule---is the critical open problem.  No
known spectral-geometric invariant (spectral determinant, Ray--Singer
analytic torsion, Seeley--DeWitt heat kernel, Atiyah--Patodi--Singer
$\eta$-invariant) reproduces $K$.  The rule mixes a discrete
combinatorial quantity ($B = 42$) with a continuous spectral invariant
($F = \pi^2/6$) in a manner without precedent in the spectral geometry
literature.

\textit{Negative result:} Applying the combination rule to the three
higher-dimensional Hopf fibrations ($S^3 \to S^7 \to S^4$,
$S^7 \to S^{15} \to S^8$) does not yield known coupling constants.
Only the complex Hopf fibration produces a physically meaningful
result~\cite{Paper2}.

%% ========================================================================
%% VII. ATOMIC SPECTRA: A TESTABLE PREDICTION PROGRAM
%% ========================================================================

\section{Atomic Spectra: A Testable Prediction Program}
\label{sec:spectra}

\subsection{Non-relativistic energies from $S^{3N-1}$}

For single-center systems, the hyperspherical framework on $S^{3N-1}$
is mathematically exact and computationally convergent.  Current
results include He at 0.004\% (2D variational + cusp correction on $S^5$,
Paper~13~\cite{Paper13}), Li at 1.10\%, and Be at 0.90\%
(multi-electron FCI).

The coupled-channel framework, already prototyped through $l_{\max} = 5$
with algebraic convergence $\sim l_{\max}^{-2}$ to a structural floor
of 0.19--0.20\%~\cite{Paper13}, provides a clear path to higher
precision.  The literature benchmark for hyperspherical He is 0.1~ppm
with 11 adiabatic channels~\cite{Lin1995}.

\subsection{Fine structure as a test of the $\alpha$ conjecture}

The non-relativistic energy $E_{\text{NR}}$ computed from $S^{3N-1}$ can
be compared with the experimental energy:
\begin{equation}
E_{\text{exp}} = E_{\text{NR}} + E_{\text{rel}}(\alpha)
+ E_{\text{QED}}(\alpha) + \cdots
\label{eq:energy_decomposition}
\end{equation}
If $E_{\text{NR}}$ is computed to sufficient precision, the residual
$E_{\text{exp}} - E_{\text{NR}}$ isolates the $\alpha$-dependent
relativistic and QED corrections.  Fine structure splitting scales as
$\alpha^2$; the Lamb shift as $\alpha^5$.  An independent determination
of $\alpha$ from geometric atomic spectra would test the Hopf conjecture
computationally.

\subsection{Milestones}

A concrete research program proceeds in four phases:
\begin{enumerate}
\item \textbf{He ground state to $<$0.01\% [ACHIEVED].}  The 2D
  variational solver on $S^5$ achieves 0.004\% with Schwartz cusp
  correction at $l_{\max} = 4$ (v2.6.0).  The adiabatic 0.20\% floor
  is bypassed by treating the hyperradius and hyperangle simultaneously.

\item \textbf{He excited states and transition energies.}  Separate
  adiabatic potential curves for ${}^1S$ and ${}^3S$ states.  Transition
  energies are more precisely measured than absolute energies; residuals
  vs.\ experiment constrain $\alpha^2$ corrections.

\item \textbf{$\alpha$ extraction.}  Compute non-relativistic He energy
  to precision better than the fine structure splitting
  ($\sim$0.36~cm$^{-1}$ for He~$2\,{}^3P$).  Test: does the extracted
  $\alpha$ match $1/137.035987$ (Hopf) vs.\ $1/137.035999$ (CODATA)?
  The $8.8 \times 10^{-8}$ difference corresponds to $\sim$0.0001~cm$^{-1}$
  in He fine structure.

\item \textbf{Li and Be.}  Ground-state energies to $< 0.01\%$ via
  coupled channels on $S^8$ (Li, 3 electrons) and $S^{11}$ (Be, 4
  electrons).  The H-set coupling from Sec.~\ref{sec:multielectron}
  provides the angular framework for Be.  Multi-electron convergence
  validates the $S^{3N-1}$ framework beyond He.
\end{enumerate}

\noindent
\textit{Status (v2.6.0).}  The 2D variational solver on~$S^5$
achieves 0.022\% raw at $l_{\max} = 7$ and 0.004\% with Schwartz
cusp correction at $l_{\max} = 4$, breaking the adiabatic 0.20\%
floor (Paper~13, Note added v2.6.0).  Phase~1 target ($< 0.01\%$)
is achieved.  The graph-native CI provides the complementary
algebraic path: 0.19\% at $n_{\max} = 7$ with zero parameters and
a fully rational FCI matrix.  Sprint~3D verified FCI basis invariance
($< 3 \times 10^{-15}$~Ha), confirming the convergence floor is
from the cusp, not basis mismatch.  Both the continuous (2D
variational) and algebraic (graph-native) paths converge to the
same exact answer from different directions.

%% ========================================================================
%% VIII. COMPUTATIONAL PIPELINE: FROM ATOMS TO MOLECULES
%% ========================================================================

\section{Computational Pipeline: From Atoms to Molecules}
\label{sec:pipeline}

The mathematical framework of
Secs.~\ref{sec:kinematics}--\ref{sec:exchange} is implemented as a
complete computational pipeline, from single atoms on $S^3$ through
polyatomic molecules via composed fiber bundles.  This section
describes the pipeline's structure, the molecular architecture, and the
30-molecule library it produces.

\subsection{Natural geometry hierarchy}

Each electron configuration has a natural coordinate system where the
Schr\"odinger equation partially or fully separates.  The choice of
geometry is determined by the symmetry of the physical system: one-center
one-electron problems retain full SO(4) symmetry on~$S^3$; each
additional center or electron breaks a symmetry and introduces a
continuous radial-like coordinate.  The hierarchy is summarized in
Table~\ref{tab:hierarchy}.

\begin{table}[t]
\caption{Natural geometry hierarchy.  Each level uses the coordinate
system where the physics naturally separates.  ``Best result'' is the
current accuracy for the prototypical system at that level.}
\label{tab:hierarchy}
\begin{ruledtabular}
\begin{tabular}{cllll}
Level & System & Geometry & Best result & Ref. \\
\hline
1 & H & $S^3$ (Fock) & $<0.1$\% & \cite{Paper7} \\
2 & H$_2^+$ & Prolate spheroid & 0.0002\% & \cite{Paper11} \\
3 & He & $S^5$ (hypersp.) & 0.004\% & \cite{Paper13} \\
4 & H$_2$ & Mol-frame hypersp. & 96.0\% $D_e$ & \cite{Paper15} \\
5 & LiH, BeH$_2$, \ldots & Composed bundles & 0.20\% $E$ & \cite{Paper17} \\
\end{tabular}
\end{ruledtabular}
\end{table}

At Level~1, Fock's SO(4) symmetry makes the entire problem discrete:
the graph Laplacian is the exact kinetic operator.  At Level~2, the
second nucleus breaks SO(4) to axial symmetry, introducing a continuous
radial coordinate in prolate spheroidal geometry~\cite{Paper11}.  At
Level~3, the second electron introduces a hyperradius~$R$ and
hyperangle~$\alpha$ on~$S^5$~\cite{Paper13}.  At Level~4, two centers
plus two electrons produce molecule-frame hyperspherical coordinates on
SO(6)~\cite{Paper15}.  In each case, the angular channels remain
discrete---labeled by quantum numbers and coupled through Gaunt
integrals---while the radial-like coordinates require numerical or
spectral solvers.

\subsection{The composed architecture}

For systems beyond H$_2$, the full $N$-electron hyperspherical space
$S^{3N-1}$ becomes computationally prohibitive (the angular basis grows
combinatorially).  The composed architecture~\cite{Paper17} factorizes
the problem by assigning each electron group---core pair, bond pair,
lone pair---its own natural coordinate system:
\begin{equation}
G_{\text{total}} = G_{\text{nuc}} \ltimes G_{\text{core}}(R)
\ltimes G_{\text{val}}(R, \text{core}).
\label{eq:composed_bundle}
\end{equation}
Each fiber is a Level~3 or Level~4 sub-problem solved independently.
The fibers couple through three mechanisms: (i)~$Z_{\text{eff}}(r)$
screening from inner shells, (ii)~Slater direct integrals $F^0$ for
inter-group Coulomb repulsion, and (iii)~a Phillips-Kleinman
pseudopotential (PK) enforcing core-valence orthogonality.

Gaunt selection rules make electron repulsion integrals (ERIs)
block-diagonal between groups: cross-block ERIs vanish identically
because the orbitals in different blocks have independent angular
quantum numbers.  This structural sparsity---not an approximation but
a selection rule---produces $O(Q^{2.5})$ Pauli term scaling in the
Jordan-Wigner qubit encoding, compared with $O(Q^{3.9\text{--}4.3})$
for Gaussian baselines~\cite{Paper14}.

\subsection{Balanced coupled composition}

The composed architecture's PK pseudopotential is the accuracy
bottleneck (5.3--26\% equilibrium distance error).  The balanced coupled
variant~\cite{Paper19} eliminates PK entirely by adding cross-center
nuclear attraction integrals: each orbital feels all nuclei, not just
its own center.

The cross-center matrix elements
$\langle \phi_a | V_B | \phi_b \rangle$ are computed via a multipole
expansion of $1/|\mathbf{r} - \mathbf{R}_B|$ in the Sturmian basis
centered on nucleus~$A$.  By the Gaunt triangle inequality, the
expansion terminates exactly at $L = 2l_{\max}$, yielding a finite,
analytically evaluated sum (incomplete gamma functions, machine
precision, zero quadrature).

The balanced coupled LiH Hamiltonian at $n_{\max} = 3$ ($Q = 84$
qubits) gives $E = -8.055$~Ha (0.20\% error vs.\ exact $-8.071$~Ha),
a $9\times$ improvement over $n_{\max} = 2$.  It is the only
configuration among four tested that produces a bound 4-electron LiH
at $n_{\max} = 2$.  Extension to BeH$_2$ (2,652~Pauli terms, 10.7\%
energy error) and H$_2$O (5,798~Pauli terms) is implemented.  The
method is optimal for single-point quantum simulation at fixed
geometries rather than potential energy surface scanning.

\subsection{The 30-molecule library}

The composed and balanced architectures have been applied to 30
molecules spanning three rows of the periodic table:
\begin{itemize}
\item \textbf{First row} ($Z = 1$--10, 8 systems): H$_2$, He, LiH,
  BeH$_2$, H$_2$O, HF, NH$_3$, CH$_4$.  Composed encoding with He
  frozen cores.
\item \textbf{Second row} ($Z = 11$--18, 6 systems): NaH, MgH$_2$,
  HCl, H$_2$S, PH$_3$, SiH$_4$.  [Ne] frozen cores via
  Clementi-Raimondi $Z_{\text{eff}}(r)$ screening.
\item \textbf{Third row} ($Z = 19$--36, 6 systems): KH, CaH$_2$,
  GeH$_4$, AsH$_3$, H$_2$Se, HBr.  [Ar] or [Ar]$3d^{10}$ frozen
  cores.
\item \textbf{Multi-center} (8 systems): LiF, CO, N$_2$, F$_2$,
  NaCl, CH$_2$O, C$_2$H$_2$, C$_2$H$_6$.
\item \textbf{Transition metals} (2 systems): ScH, TiH.  [Ar] frozen
  cores with $d$-orbital valence blocks.
\end{itemize}

A universal linear scaling law holds across the entire library:
\begin{equation}
N_{\text{Pauli}} = 11.11 \times Q \quad (\pm\, 0.1),
\label{eq:linear_scaling}
\end{equation}
where $Q$ is the qubit count.  This coefficient is stable across all
30 molecules ($Q = 10$--160), independent of block count, molecular
structure, or atomic species.  Isostructural invariance is exact:
molecules with identical block topology produce identical Pauli counts
regardless of nuclear charges (NaH $=$ KH $= 239$;
CO $=$ N$_2$ $=$ F$_2$ $= 1{,}111$).  The $d$-orbital blocks in
transition metal hydrides are sparser still, with Pauli/$Q = 9.27$
(ERI density 4.0\% vs.\ 8.9\% for $s$+$p$ blocks).

\subsection{Code access}

The composed Hamiltonians for all 30 molecules are available as a
standalone Python package:
\begin{verbatim}
pip install geovac-hamiltonians
from geovac.ecosystem_export import hamiltonian
h = hamiltonian("LiH", R=3.015)
qop = h.to_openfermion()
\end{verbatim}
The \texttt{hamiltonian()} API returns qubit operators in OpenFermion,
Qiskit, and PennyLane formats.  The package is a 92~KB wheel with no
heavy dependencies.  The full source code, including all solvers,
test suite ($\sim$1000 tests), and paper sources, is available at
\texttt{github.com/jloutey-hash/geovac} with Zenodo DOI-stamped
releases~\cite{Paper20}.

%% ========================================================================
%% IX. SCOPE AND HONEST LIMITATIONS
%% ========================================================================

\section{Scope and Honest Limitations}
\label{sec:limitations}

\subsection{What the framework does}

\begin{enumerate}
\item Derives quantum number structure $(l, m, 2l+1)$ from geometric
  axioms, independent of the Schr\"odinger equation.
\item Provides a mathematically exact formulation for single-center
  Coulomb systems, verified through 18 symbolic proofs.
\item Produces computationally efficient qubit encodings:
  $O(Q^{2.5})$ Pauli scaling for composed molecular Hamiltonians,
  51--1,712$\times$ fewer Pauli terms than Gaussian
  baselines~\cite{Paper14}.
\item Identifies a new sparsity mechanism (6$j$ recoupling) that is
  basis-agnostic.
\item Classifies transcendental content through the exchange constant
  taxonomy.
\item Suggests a topological origin for $\alpha$ with zero free
  parameters and $p = 5.2 \times 10^{-9}$ (conjectural).
\end{enumerate}

\subsection{What the framework does not do}

\begin{enumerate}
\item Replace computational quantum chemistry for molecules.  The
  molecular limitation is fundamental: the multi-center $N$-body
  problem does not separate in any known coordinate system.
\item Provide convergent molecular potential energy surfaces.  The
  composed approximation has characterized accuracy ceilings
  (5.3--26\% equilibrium distance error).
\item Derive the combination rule for $\alpha$ (Link~3 remains open).
\item Account for relativistic effects.  The framework is
  non-relativistic; fine structure enters as a correction
  (Sec.~\ref{sec:spectra}).
\item Explain why $S^3$ is the physical manifold.  The framework takes
  Fock's 1935 result as given and builds upon it.
\end{enumerate}

\subsection{The molecular boundary}

Six nested approaches to molecular electronic structure within the
$S^{3N-1}$ framework have been tested: single-center, charge-center,
heterogeneous Löwdin, and three variants of the composed
architecture.  All single-basis approaches fail for heteronuclear
molecular potential energy surfaces.  The root cause is that core and
valence electrons occupy different length scales; any unified-basis
approach faces orthogonalization penalties that destroy the sparsity
advantage.  The composed architecture~\cite{Paper17}---which assigns
each electron group its own natural coordinate system---provides
practical molecular Hamiltonians with $N_{\text{Pauli}} \approx
11.11 \times Q$ across 30 molecules.

This molecular limitation does not affect the single-center mathematical
structure, just as turbulence does not invalidate the Navier--Stokes
equations.

%% ========================================================================
%% IX. CONCLUSION
%% ========================================================================

\section{Conclusion}
\label{sec:conclusion}

The kinematic structure of non-relativistic quantum
mechanics---quantum numbers, degeneracies, selection rules---emerges
from the spectral theory of Laplace--Beltrami operators on compact
manifolds, specifically the unit hyperspheres $S^{3N-1}$ for
$N$-electron systems.  The dynamic structure---energy levels, the
Coulomb potential---enters through projection parameters that map the
dimensionless topology into physical observables.  Transcendental
content is classified by the exchange constant taxonomy, appearing at
precisely identifiable steps in the projection chain.

For single-center systems, these results are mathematically proven and
computationally validated: 18 symbolic proofs for $S^3$, a 0.004\% He
ground state on $S^5$ via 2D variational solution with Schwartz cusp
correction~\cite{Paper13}, chemical periodicity from $S_N$
representation theory on $S^{3N-1}$~\cite{Paper16}, and a 6$j$
recoupling sparsity mechanism reducing ERI density by 18--26\%.  For
molecules, the composed architecture produces structurally sparse qubit
Hamiltonians with $N_{\text{Pauli}} \approx 11.11 \times Q$ across 30
molecules ($Z = 1$--36), $O(Q^{2.5})$ Pauli scaling, and balanced
coupled FCI energy at 0.20\% error for
LiH~\cite{Paper17,Paper19,Paper20}.

The Hopf fibration of the electron's $S^3$ state manifold suggests---but
does not yet prove---that the electromagnetic coupling constant $\alpha$
is determined by the same topology~\cite{Paper2}.  A systematic program
of high-precision atomic spectra computation on $S^{3N-1}$, compared
with experimental fine structure data, offers a concrete and testable
path forward.

The full framework is available as open-source software
(\texttt{pip install geovac-hamiltonians}), with ecosystem export to
OpenFermion, Qiskit, and PennyLane~\cite{Paper20}.

%% ========================================================================
%% ACKNOWLEDGMENTS
%% ========================================================================

\begin{acknowledgments}
This work was developed using an AI-augmented agentic workflow, with
implementation and documentation drafting performed collaboratively
with Anthropic's Claude.
\end{acknowledgments}

%% ========================================================================
%% REFERENCES
%% ========================================================================

\begin{thebibliography}{99}

\bibitem{Fock1935}
V.~Fock, Z.\ Phys.\ \textbf{98}, 145 (1935).

\bibitem{Bargmann1936}
V.~Bargmann, Z.\ Phys.\ \textbf{99}, 576 (1936).

\bibitem{BarutKleinert1967}
A.~O.~Barut and H.~Kleinert, Phys.\ Rev.\ \textbf{156}, 1541 (1967).

\bibitem{Hopf1931}
H.~Hopf, Math.\ Ann.\ \textbf{104}, 637 (1931).

\bibitem{Macek1968}
J.~Macek, J.\ Phys.\ B \textbf{1}, 831 (1968).

\bibitem{Lin1995}
C.~D.~Lin, Phys.\ Rep.\ \textbf{257}, 1 (1995).

\bibitem{Paper0}
J.~Loutey, ``Geometric packing and the origin of quantum numbers,''
GeoVac Paper~0 (2025).

\bibitem{Paper2}
J.~Loutey, ``The fine structure constant from Hopf bundle spectral
invariants,'' GeoVac Paper~2 (2025).

\bibitem{Paper7}
J.~Loutey, ``The dimensionless vacuum: Schr\"odinger recovery from the
unit $S^3$,'' GeoVac Paper~7 (2025).

\bibitem{Paper13}
J.~Loutey, ``Hyperspherical lattice for two-electron atoms,'' GeoVac
Paper~13 (2026).

\bibitem{Paper14}
J.~Loutey, ``Structurally sparse qubit Hamiltonians from natural
geometry,'' GeoVac Paper~14 (2026).

\bibitem{Paper15}
J.~Loutey, ``Molecule-frame hyperspherical geometry for $\text{H}_2$,''
GeoVac Paper~15 (2026).

\bibitem{Paper16}
J.~Loutey, ``Chemical periodicity as $S_N$ representation theory,''
GeoVac Paper~16 (2026).

\bibitem{Paper17}
J.~Loutey, ``Composed natural geometries for polyatomic molecules,''
GeoVac Paper~17 (2026).

\bibitem{Paper18}
J.~Loutey, ``Spectral-geometric exchange constants: Projecting discrete
spectra onto continuous manifolds,'' GeoVac Paper~18 (2026).

\bibitem{Paper11}
J.~Loutey, ``The prolate spheroidal lattice for H$_2^+$,'' GeoVac
Paper~11 (2026).

\bibitem{Paper19}
J.~Loutey, ``Balanced coupled composition: PK-free molecular
Hamiltonians via Shibuya-Wulfman integrals,'' GeoVac Paper~19 (2026).

\bibitem{Paper20}
J.~Loutey, ``Quantum resource benchmarks for structurally sparse
Hamiltonians,'' GeoVac Paper~20 (2026).

\end{thebibliography}

\end{document}
